\chapter{2023年全国甲卷（理）}

\section{单选题}

\begin{problem}
设 \( U=\Z \)，集合 \( M=\{x\mid x=3k+1,\ k\in \Z\} \)，\( N=\{x\mid x=3k+2,\ k\in \Z\} \)，则 \( \complement_U\left( M\cup N \right) \) 为（\quad）

\choices
    {\(\{x\mid x=3k,\ k\in \Z\}\)}
    {\(\{x\mid x=3k-1,\ k\in \Z\}\)}
    {\(\{x\mid x=3k-2,\ k\in \Z\}\)}
    {\(\varnothing\)}

\end{problem}

\begin{answer}
A
\end{answer}

\begin{solution}
因为
\(\Z=\{x\mid x=3k,\ k\in \Z\}\cup\{x\mid x=3k+1,\ k\in \Z\}\cup\{x\mid x=3k+2,\ k\in \Z\}\)，
且 \( U=\Z \)，

所以
\(\complement_U\left( M\cup N \right)=\{x\mid x=3k,\ k\in \Z\}\).

故选 \( A \).
\end{solution}

\begin{problem}
设 \( a\in \R \)，\( (a+\i)(1-\i a)=2 \)，则 \( a= \)（\quad）

\choices
    {\(-1\)}
    {0}
    {1}
    {2}

\end{problem}

\begin{answer}
C
\end{answer}

\begin{solution}
因为
\[
(a+\i)(1-\i a)=a-a^2\i+\i+a=2a+\left( 1-a^2 \right)\i=2
\]
所以
\(\begin{cases}
2a=2,\\
1-a^2=0,
\end{cases}\)
解得 \( a=1 \).

故选 \( C \).
\end{solution}

\begin{problem}
执行下面的程序框图，输出的 \( B= \)（\quad）

\choices
    {21}
    {34}
    {55}
    {89}

\end{problem}

\begin{answer}
B
\end{answer}

\begin{solution}
当 \( k=1 \) 时，判断框条件满足，第一次执行循环体，\( A=1+2=3 \)，\( B=3+2=5 \)，\( k=1+1=2 \)；

当 \( k=2 \) 时，判断框条件满足，第二次执行循环体，\( A=3+5=8 \)，\( B=8+5=13 \)，\( k=2+1=3 \)；

当 \( k=3 \) 时，判断框条件满足，第三次执行循环体，\( A=8+13=21 \)，\( B=21+13=34 \)，\( k=3+1=4 \)；

当 \( k=4 \) 时，判断框条件不满足，跳出循环体，输出 \( B=34 \).

故选 \( B \).
\end{solution}

\begin{problem}
已知向量 \( \bs{a},\bs{b},\bs{c} \) 满足 \( \lvert \bs{a} \rvert=\lvert \bs{b} \rvert=1 \)，\( \lvert \bs{c} \rvert=\sqrt{2} \)，且 \( \bs{a}+\bs{b}+\bs{c}=\bs{0} \)，则 \( \cos\left( \bs{a}-\bs{c},\bs{b}-\bs{c} \right)= \)（\quad）

\bitmapfigure[max width=6.0cm,max totalheight=4.8cm]{img/2023/national_paper_A_science/q4_fig1.png}

\choices
    {\(-\frac45\)}
    {\(-\frac25\)}
    {\(\frac25\)}
    {\(\frac45\)}

\end{problem}

\begin{answer}
D
\end{answer}

\begin{solution}
因为 \( \bs{a}+\bs{b}+\bs{c}=\bs{0} \)，所以 \( \bs{a}+\bs{b}=-\bs{c} \)，

即
\(\lvert \bs{a} \rvert^2+\lvert \bs{b} \rvert^2+2\bs{a}\cdot\bs{b}=\lvert \bs{c} \rvert^2\)，
所以
\(1+1+2\bs{a}\cdot\bs{b}=2\)，
从而 \( \bs{a}\cdot\bs{b}=0 \).

如图，设 \( \overrightarrow{OA}=\bs{a} \)，\( \overrightarrow{OB}=\bs{b} \)，\( \overrightarrow{OC}=\bs{c} \). 由题知 \( OA=OB=1 \)，\( OC=\sqrt{2} \)，故 \( \triangle OAB \) 是等腰直角三角形.

设 \( D \) 为 \( AB \) 的中点，则
\(OD=\frac{\sqrt{2}}{2}\)，\(AD=\frac{\sqrt{2}}{2}\)，
从而
\(CD=CO+OD=\sqrt{2}+\frac{\sqrt{2}}{2}=\frac{3\sqrt{2}}{2}\).

于是
\(\tan\angle ACD=\frac{AD}{CD}=\frac13\)，
所以
\(\cos\angle ACD=\frac{3}{\sqrt{10}}\).

故
\[
\cos\left( \bs{a}-\bs{c},\bs{b}-\bs{c} \right)=\cos\angle ACB=\cos2\angle ACD=2\cos^2\angle ACD-1=2\times\left( \frac{3}{\sqrt{10}} \right)^2-1=\frac45
\]

故选 \( D \).
\end{solution}

\begin{problem}
设等比数列 \( \{a_n\} \) 的各项均为正数，前 \( n \) 项和为 \( S_n \)，若 \( a_1=1 \)，\( S_5=5S_3-4 \)，则 \( S_4= \)（\quad）

\choices
    {\(\frac{15}{8}\)}
    {\(\frac{65}{8}\)}
    {15}
    {40}

\end{problem}

\begin{answer}
C
\end{answer}

\begin{solution}
设公比为 \( q \). 由题知
\(1+q+q^2+q^3+q^4=5\left( 1+q+q^2 \right)-4\)，
即
\(q^3+q^4=4q+4q^2\)，
所以
\(q^3+q^2-4q-4=0\)，
即
\(\left( q-2 \right)\left( q+1 \right)\left( q+2 \right)=0\).

由题知 \( q>0 \)，所以 \( q=2 \).

因此
\(S_4=1+2+4+8=15\).

故选 \( C \).
\end{solution}

\begin{problem}
某地的中学生中有 60\% 的同学爱好滑冰，50\% 的同学爱好滑雪，70\% 的同学爱好滑冰或爱好滑雪. 在该地的中学生中随机调查一位同学，若该同学爱好滑雪，则该同学也爱好滑冰的概率为（\quad）

\choices
    {0.8}
    {0.6}
    {0.5}
    {0.4}

\end{problem}

\begin{answer}
A
\end{answer}

\begin{solution}
同时爱好两项的概率为 \( 0.5+0.6-0.7=0.4 \).

记“该同学爱好滑雪”为事件 \( A \)，记“该同学爱好滑冰”为事件 \( B \)，则 \( P(A)=0.5 \)，\( P(AB)=0.4 \).

所以
\(P(B\mid A)=\frac{P(AB)}{P(A)}=\frac{0.4}{0.5}=0.8\).

故选 \( A \).
\end{solution}

\begin{problem}
设甲：\( \sin^2\alpha+\sin^2\beta=1 \)，乙：\( \sin\alpha+\cos\beta=0 \). 则（\quad）

\choices
    {甲是乙的充分条件但不是必要条件}
    {甲是乙的必要条件但不是充分条件}
    {甲是乙的充要条件}
    {甲既不是乙的充分条件也不是乙的必要条件}

\end{problem}

\begin{answer}
B
\end{answer}

\begin{solution}
当 \( \sin^2\alpha+\sin^2\beta=1 \) 时，例如 \( \alpha=\frac{\pi}{2} \)，\( \beta=0 \)，但 \( \sin\alpha+\cos\beta\neq 0 \)，

即甲推不出乙.

当 \( \sin\alpha+\cos\beta=0 \) 时，
\(\sin^2\alpha+\sin^2\beta=\left( -\cos\beta \right)^2+\sin^2\beta=1\)，
即乙能推出甲.

综上可知，甲是乙的必要不充分条件.

故选 \( B \).
\end{solution}

\begin{problem}
已知双曲线 \( C:\frac{x^2}{a^2}-\frac{y^2}{b^2}=1 \)（\(a>0\)，\(b>0\)）的离心率为 \( \sqrt{5} \)，\( C \) 的一条渐近线与圆 \( \left( x-2 \right)^2+\left( y-3 \right)^2=1 \) 交于 \( A,B \) 两点，则 \( \lvert AB \rvert= \)（\quad）

\choices
    {\(\frac{\sqrt{5}}{5}\)}
    {\(\frac{2\sqrt{5}}{5}\)}
    {\(\frac{3\sqrt{5}}{5}\)}
    {\(\frac{4\sqrt{5}}{5}\)}

\end{problem}

\begin{answer}
D
\end{answer}

\begin{solution}
由 \( e=\sqrt{5} \)，则
\[
\frac{c^2}{a^2}=\frac{a^2+b^2}{a^2}=1+\frac{b^2}{a^2}=5
\]
解得
\(\frac{b}{a}=2\).

所以双曲线的一条渐近线不妨取 \( y=2x \). 则圆心 \( (2,3) \) 到渐近线的距离为
\(d=\frac{\lvert 2\times 2-3 \rvert}{\sqrt{2^2+1}}=\frac{\sqrt{5}}{5}\).

因此弦长
\(\lvert AB \rvert=2\sqrt{r^2-d^2}=2\sqrt{1-\frac15}=\frac{4\sqrt{5}}{5}\).

故选 \( D \).
\end{solution}

\begin{problem}
现有 5 名志愿者报名参加公益活动，在某一星期的星期六、星期日两天，每天从这 5 人中安排 2 人参加公益活动，则恰有 1 人在这两天都参加的不同安排方式共有（\quad）

\choices
    {120}
    {60}
    {30}
    {20}

\end{problem}

\begin{answer}
B
\end{answer}

\begin{solution}
不妨记五名志愿者为 \( a,b,c,d,e \).

先从 \( a,b,c,d,e \) 中选出 1 人，作为连续两天都参加公益活动的人，有 5 种选法.

连续参加两天公益活动的人确定后，再从剩余的 4 人中抽取 2 人，分别安排参加星期六与星期天的公益活动，共有 \( A_4^2=12 \) 种方法.

所以恰有 1 人连续参加了两天公益活动的选择种数有
\(5\times 12=60\).

故选 \( B \).
\end{solution}

\begin{problem}
函数 \( y=f(x) \) 的图象由函数 \( y=\cos\left( 2x+\frac{\pi}{6} \right) \) 的图象向左平移 \( \frac{\pi}{6} \) 个单位长度得到，则 \( y=f(x) \) 的图象与直线 \( y=\frac12x-\frac12 \) 的交点个数为（\quad）

\choices
    {1}
    {2}
    {3}
    {4}

\end{problem}

\begin{answer}
C
\end{answer}

\begin{solution}
因为函数 \( y=\cos\left( 2x+\frac{\pi}{6} \right) \) 向左平移 \( \frac{\pi}{6} \) 个单位所得函数为 \( y=\cos\left[ 2\left( x+\frac{\pi}{6} \right)+\frac{\pi}{6} \right]=\cos\left( 2x+\frac{\pi}{2} \right)=-\sin2x \)，所以 \( f(x)=-\sin2x \).

当 \( x=0 \) 时，\( y=\frac12x-\frac12=-\frac12 \)；当 \( x=1 \) 时，\( y=\frac12x-\frac12=0 \). 所以直线 \( y=\frac12x-\frac12 \) 过 \( \left( 0,-\frac12 \right) \) 与 \( (1,0) \) 两点.

考虑 \(2x=-\frac{3\pi}{2}\)，\(2x=\frac{3\pi}{2}\)，\(2x=\frac{7\pi}{2}\)，即 \(x=-\frac{3\pi}{4}\)，\(x=\frac{3\pi}{4}\)，\(x=\frac{7\pi}{4}\) 处 \( f(x) \) 与 \( y=\frac12x-\frac12 \) 的大小关系：
当 \(x=-\frac{3\pi}{4}\) 时，\(f\left( -\frac{3\pi}{4} \right)=-\sin\left( -\frac{3\pi}{2} \right)=-1\)，且 \(\frac12\times\left( -\frac{3\pi}{4} \right)-\frac12=\frac{-3\pi-4}{8}<-1\)；当 \(x=\frac{3\pi}{4}\) 时，\(f\left( \frac{3\pi}{4} \right)=-\sin\frac{3\pi}{2}=1\)，且 \(\frac12\times\frac{3\pi}{4}-\frac12=\frac{3\pi-4}{8}<1\)；当 \(x=\frac{7\pi}{4}\) 时，\(f\left( \frac{7\pi}{4} \right)=-\sin\frac{7\pi}{2}=1\)，且 \(\frac12\times\frac{7\pi}{4}-\frac12=\frac{7\pi-4}{8}>1\).

\bitmapfigure[max width=5.2cm,max totalheight=4.8cm]{img/2023/national_paper_A_science/q10_solution_fig1.png}

所以由图可知，\( f(x) \) 与 \( y=\frac12x-\frac12 \) 的交点个数为 3.

故选 \( C \).
\end{solution}

\begin{problem}
已知四棱锥 \( P-ABCD \) 的底面是边长为 4 的正方形，\( PC=PD=3 \)，\( \angle PCA=45^\circ \)，则 \( \triangle PBC \) 的面积为（\quad）

\choices
    {\(2\sqrt{2}\)}
    {\(3\sqrt{2}\)}
    {\(4\sqrt{2}\)}
    {\(6\sqrt{2}\)}

\end{problem}

\begin{answer}
C
\end{answer}

\begin{solution}
连结 \( AC,BD \) 交于 \( O \)，连结 \( PO \)，则 \( O \) 为 \( AC,BD \) 的中点.

因为底面 \( ABCD \) 为正方形，\( AB=4 \)，所以
\(AC=BD=4\sqrt{2}\).

又 \( PC=PD=3 \)，\( DO=CO=2\sqrt{2} \)，\( PO=OP \)，所以 \( \triangle PDO\cong \triangle PCO \)，从而 \( \angle PDO=\angle PCO \).

又 \( PC=PD \)，\( AC=BD \)，所以 \( \triangle PDB\cong \triangle PCA \)，故 \( PA=PB \).

在 \( \triangle PAC \) 中，\( PC=3 \)，\( AC=4\sqrt{2} \)，\( \angle PCA=45^\circ \)，
\[
PA^2=AC^2+PC^2-2AC\cdot PC\cos\angle PCA=32+9-2\times 4\sqrt{2}\times 3\times \frac{\sqrt{2}}{2}=17
\]
所以 \( PA=\sqrt{17} \)，从而 \( PB=\sqrt{17} \).

于是，在 \( \triangle PBC \) 中，
\(PC=3\)，\(PB=\sqrt{17}\)，\(BC=4\)，
故
\(\cos\angle PCB=\frac{PC^2+BC^2-PB^2}{2PC\cdot BC}=\frac{9+16-17}{2\times 3\times 4}=\frac13\).

又 \( 0<\angle PCB<\pi \)，所以 \(\sin\angle PCB=\sqrt{1-\cos^2\angle PCB}=\sqrt{1-\frac19}=\frac{2\sqrt{2}}{3}\).

因此 \( S_{\triangle PBC}=\frac12PC\cdot BC\cdot \sin\angle PCB=\frac12\times 3\times 4\times \frac{2\sqrt{2}}{3}=4\sqrt{2} \).

故选 \( C \).
\end{solution}

\begin{problem}
设 \( O \) 为坐标原点，\( F_1,F_2 \) 为椭圆 \( \frac{x^2}{9}+\frac{y^2}{6}=1 \) 的两个焦点，点 \( P \) 在该椭圆上，且 \( \cos\angle F_1PF_2=\frac35 \). 则 \( \lvert OP \rvert= \)（\quad）

\choices
    {\(\frac{13}{5}\)}
    {\(\frac{\sqrt{30}}{2}\)}
    {\(\frac{14}{5}\)}
    {\(\frac{\sqrt{35}}{2}\)}

\end{problem}

\begin{answer}
B
\end{answer}

\begin{solution}
设 \( \angle F_1PF_2=2\theta \).

因为
\(\cos2\theta=\frac{1-\tan^2\theta}{1+\tan^2\theta}=\frac35\)，
解得
\(\tan\theta=\frac12\).

由椭圆方程可知
\(a^2=9\)，\(b^2=6\)，\(c^2=a^2-b^2=3\).

由焦点三角形面积公式得
\(S_{\triangle PF_1F_2}=b^2\tan\theta=6\times \frac12=3\).

另一方面，\( S_{\triangle PF_1F_2}=\frac12\cdot F_1F_2\cdot \lvert y_P \rvert=\frac12\cdot 2\sqrt{3}\cdot \lvert y_P \rvert=\sqrt{3}\lvert y_P \rvert \)，所以 \( \lvert y_P \rvert=\sqrt{3} \)，即 \( y_P^2=3 \).

代入椭圆方程得
\(\frac{x_P^2}{9}+\frac36=1\)，
所以
\(x_P^2=\frac92\).

因此 \( \lvert OP \rvert=\sqrt{x_P^2+y_P^2}=\sqrt{\frac92+3}=\sqrt{\frac{15}{2}}=\frac{\sqrt{30}}{2} \).

故选 \( B \).
\end{solution}
\section{填空题}

\begin{problem}
设 \( f(x)=\left( x-1 \right)^2+ax+\sin\left( x+\frac{\pi}{2} \right) \). 若 \( f(x) \) 为偶函数，则 \( a= \)（\fillinblank{}）.

\end{problem}

\begin{answer}
2
\end{answer}

\begin{solution}
因为
\(f(x)=\left( x-1 \right)^2+ax+\sin\left( x+\frac{\pi}{2} \right)=\left( x-1 \right)^2+ax+\cos x\)
为偶函数，定义域为 \( \R \)，

所以 \(f\left( -\frac{\pi}{2} \right)=f\left( \frac{\pi}{2} \right)\)，即 \(\left( -\frac{\pi}{2}-1 \right)^2-\frac{\pi}{2}a+\cos\left( -\frac{\pi}{2} \right)=\left( \frac{\pi}{2}-1 \right)^2+\frac{\pi}{2}a+\cos\frac{\pi}{2}\).

于是
\(\pi a=\left( \frac{\pi}{2}+1 \right)^2-\left( \frac{\pi}{2}-1 \right)^2=2\pi\)，
故 \( a=2 \).

此时
\(f(x)=\left( x-1 \right)^2+2x+\cos x=x^2+1+\cos x\)，
所以
\(f(-x)=(-x)^2+1+\cos(-x)=x^2+1+\cos x=f(x)\).

故答案为：2.
\end{solution}

\begin{problem}
若 \( x,y \) 满足约束条件 \( 3x-2y\le 3 \)，\( -2x+3y\le 3 \)，\( x+y\ge 1 \)，设 \( z=3x+2y \) 的最大值为（\fillinblank{}）.

\end{problem}

\begin{answer}
15
\end{answer}

\begin{solution}
作出可行域，由线性规划可知，当目标函数
\(y=-\frac32x+\frac{z}{2}\)
过点 \( A \) 时，\( z \) 有最大值.

由
\(\begin{cases}
-2x+3y=3,\\
3x-2y=3,
\end{cases}\)
可得
\(\begin{cases}
x=3,\\
y=3,
\end{cases}\)
即 \( A(3,3) \).

所以
\(z_{\max}=3\times 3+2\times 3=15\).

故答案为：15.
\end{solution}

\begin{problem}
在正方体 \( ABCD-A_1B_1C_1D_1 \) 中，\( E,F \) 分别为 \( AB,C_1D_1 \) 的中点，以 \( EF \) 为直径的球的球面与该正方体的棱共有（\fillinblank{}）个公共点.

\end{problem}

\begin{answer}
12
\end{answer}

\begin{solution}
不妨设正方体棱长为 2，\( EF \) 中点为 \( O \)，取 \( CD,CC_1 \) 中点 \( G,M \)，侧面 \( BB_1C_1C \) 的中心为 \( N \)，连接 \( FG,EG,OM,ON,MN \).

由题意可知，\( O \) 为球心. 因为 \( E,F \) 分别为 \( AB,C_1D_1 \) 的中点，所以 \( O \) 是正方体的中心. 在正方体中，
\(EF=\sqrt{FG^2+EG^2}=\sqrt{2^2+2^2}=2\sqrt{2}\)，
即球半径
\(R=\sqrt{2}\).

则球心 \( O \) 到棱 \( CC_1 \) 的距离为
\(OM=\sqrt{ON^2+MN^2}=\sqrt{1^2+1^2}=\sqrt{2}\).

所以球与棱 \( CC_1 \) 相切，球面与棱 \( CC_1 \) 只有 1 个交点.

又因为 \( O \) 是正方体的中心，正方体的 12 条棱关于中心的位置完全对称，所以球心 \( O \) 到每条棱的距离都等于 \( OM=\sqrt{2} \). 因此球与其余各棱也都相切，球面与每条棱都只有 1 个交点.

所以以 \( EF \) 为直径的球面与正方体每条棱的交点总数为 12.

故答案为：12.
\end{solution}

\begin{problem}
在 \( \triangle ABC \) 中，\( \angle BAC=60^\circ \)，\( AB=2 \)，\( BC=\sqrt{6} \)，\( \angle BAC \) 的角平分线交 \( BC \) 于 \( D \)，则 \( AD= \)（\fillinblank{}）.

\end{problem}

\begin{answer}
2
\end{answer}

\begin{solution}
记 \( AB=c \)，\( AC=b \)，\( BC=a \).

由余弦定理可得
\(2^2+b^2-2\times 2\times b\times \cos 60^\circ=6\).
因为 \( b>0 \)，解得
\(b=1+\sqrt{3}\).

由 \(S_{\triangle ABC}=S_{\triangle ABD}+S_{\triangle ACD}\) 可得 \(\frac12\times 2\times b\times \sin 60^\circ=\frac12\times 2\times AD\times \sin 30^\circ+\frac12\times AD\times b\times \sin 30^\circ\).

解得
\(AD=\frac{\sqrt{3}\,b}{1+\frac{b}{2}}=\frac{2\sqrt{3}\left( 1+\sqrt{3} \right)}{3+\sqrt{3}}=2\).

故答案为：2.
\end{solution}
\section{解答题}

\begin{problem}
设 \( S_n \) 为数列 \( \{a_n\} \) 的前 \( n \) 项和，已知 \( a_2=1 \)，\( 2S_n=na_n \).

\begin{enumerate}
\item 求 \( \{a_n\} \) 的通项公式；
\item 求数列 \( \left\{ \frac{a_{n+1}}{2^n} \right\} \) 的前 \( n \) 项和 \( T_n \).
\end{enumerate}

\end{problem}

\begin{solution}
（1）因为 \( 2S_n=na_n \)，

当 \( n=1 \) 时，\( 2a_1=a_1 \)，即 \( a_1=0 \)；

当 \( n=3 \) 时，\( 2\left( 1+a_3 \right)=3a_3 \)，即 \( a_3=2 \).

当 \( n\ge 2 \) 时，
\(2S_{n-1}=(n-1)a_{n-1}\)，
所以
\(2\left( S_n-S_{n-1} \right)=na_n-(n-1)a_{n-1}=2a_n\).

化简得
\(\left( n-2 \right)a_n=\left( n-1 \right)a_{n-1}\).

当 \( n\ge 3 \) 时，
\(\frac{a_n}{n-1}=\frac{a_{n-1}}{n-2}=\cdots=\frac{a_3}{2}=1\)，
即
\(a_n=n-1\).

当 \( n=1,2,3 \) 时都满足上式，所以
\(a_n=n-1\)（\(n\ge 1\)）.

（2）因为
\(\frac{a_{n+1}}{2^n}=\frac{n}{2^n}\)，
所以
\(T_n=1\times\left( \frac12 \right)^1+2\times\left( \frac12 \right)^2+3\times\left( \frac12 \right)^3+\cdots+n\times\left( \frac12 \right)^n\)，
\(\frac12T_n=1\times\left( \frac12 \right)^2+2\times\left( \frac12 \right)^3+\cdots+\left( n-1 \right)\times\left( \frac12 \right)^n+n\times\left( \frac12 \right)^{n+1}\).

两式相减得
\(\frac12T_n=\left( \frac12 \right)^1+\left( \frac12 \right)^2+\left( \frac12 \right)^3+\cdots+\left( \frac12 \right)^n-n\times\left( \frac12 \right)^{n+1}\).

因此
\[
\frac12T_n=\frac{\frac12\left[ 1-\left( \frac12 \right)^n \right]}{1-\frac12}-n\times\left( \frac12 \right)^{n+1}=1-\left( 1+\frac{n}{2} \right)\left( \frac12 \right)^n
\]
故 \( T_n=2-\left( 2+n \right)\left( \frac12 \right)^n \).
\end{solution}

\begin{problem}
如图，在三棱柱 \( ABC-A_1B_1C_1 \) 中，\( A_1C\perp \) 底面 \( ABC \)，\( \angle ACB=90^\circ \)，\( AA_1=2 \)，\( A_1 \) 到平面 \( BCC_1B_1 \) 的距离为 1.

\begin{enumerate}
\item 证明 \( A_1C=AC \)；
\item 已知 \( AA_1 \) 与 \( BB_1 \) 的距离为 2，求 \( AB_1 \) 与平面 \( BCC_1B_1 \) 所成角的正弦值.
\end{enumerate}

\bitmapfigure[max width=6.0cm,max totalheight=4.8cm]{img/2023/national_paper_A_science/q18_fig1.png}

\end{problem}

\begin{solution}
（1）如图，

\bitmapfigure[max width=4.8cm,max totalheight=4.8cm]{img/2023/national_paper_A_science/q18_solution_fig1.png}

因为 \( A_1C\perp \) 底面 \( ABC \)，\( BC\subset \) 平面 \( ABC \)，所以 \( A_1C\perp BC \).

又 \( BC\perp AC \)，且 \( A_1C,AC\subset \) 平面 \( ACC_1A_1 \)，\( A_1C\cap AC=C \)，所以 \( BC\perp \text{平面 }ACC_1A_1 \).

又 \( BC\subset \) 平面 \( BCC_1B_1 \)，故
\(\text{平面 }ACC_1A_1\perp \text{平面 }BCC_1B_1\).

过 \( A_1 \) 作 \( A_1O\perp CC_1 \) 交 \( CC_1 \) 于 \( O \). 又平面 \( ACC_1A_1\cap \) 平面 \( BCC_1B_1=CC_1 \)，且 \( A_1O\subset \) 平面 \( ACC_1A_1 \)，故
\(A_1O\perp \text{平面 }BCC_1B_1\).

因为 \( A_1 \) 到平面 \( BCC_1B_1 \) 的距离为 1，所以
\(A_1O=1\).

在直角三角形 \( A_1CC_1 \) 中，\( A_1C\perp A_1C_1 \)，且
\(CC_1=AA_1=2\).

设 \( CO=x \)，则 \( C_1O=2-x \).

由
\(CO^2+A_1O^2=A_1C^2\)，
\(A_1O^2+OC_1^2=A_1C_1^2\)，
\(A_1C^2+A_1C_1^2=CC_1^2\)
可得
\(1+x^2+1+\left( 2-x \right)^2=4\)，
解得 \( x=1 \).

所以
\(AC=A_1C=A_1C_1=\sqrt{2}\)，
故 \( A_1C=AC \).

（2）由（1）知 \( AC=A_1C \)，且 \( BC\perp A_1C \)，\( BC\perp AC \)，故 \( \triangle ACB\cong \triangle A_1CB \)，于是
\(BA=BA_1\).

过 \( B \) 作 \( BD\perp AA_1 \)，交 \( AA_1 \) 于 \( D \)，则 \( D \) 为 \( AA_1 \) 中点.

由直线 \( AA_1 \) 与 \( BB_1 \) 的距离为 2，所以
\(BD=2\).

又 \( A_1D=1 \)，故
\(A_1B=AB=\sqrt{5}\).

在直角三角形 \( ABC \) 中，
\(BC=\sqrt{AB^2-AC^2}=\sqrt{5-2}=\sqrt{3}\).

延长 \( AC \)，使 \( AC=CM \)，连接 \( C_1M \).

由 \( CM//A_1C_1 \)，\( CM=A_1C_1 \) 知四边形 \( A_1CMC_1 \) 为平行四边形，所以 \( C_1M//A_1C \)，从而 \( C_1M\perp \) 平面 \( ABC \).

又 \( AM\subset \) 平面 \( ABC \)，所以 \( C_1M\perp AM \).

在直角三角形 \( AC_1M \) 中，
\(AM=2AC\)，\(C_1M=A_1C\)，
所以
\(AC_1=\sqrt{\left( 2AC \right)^2+A_1C^2}\).

在直角三角形 \( AB_1C_1 \) 中，\( B_1C_1=BC=\sqrt{3} \)，故
\(AB_1=\sqrt{\left( 2\sqrt{2} \right)^2+\left( \sqrt{2} \right)^2+\left( \sqrt{3} \right)^2}=\sqrt{13}\).

又点 \( A \) 到平面 \( BCC_1B_1 \) 的距离也为 1，所以 \( AB_1 \) 与平面 \( BCC_1B_1 \) 所成角的正弦值为 \( \frac{1}{AB_1}=\frac{1}{\sqrt{13}}=\frac{\sqrt{13}}{13} \).
\end{solution}

\begin{problem}
一项试验旨在研究臭氧效应. 实验方案如下：选 40 只小白鼠，随机地将其中 20 只分配到实验组，另外 20 只分配到对照组，实验组的小白鼠饲养在高浓度臭氧环境，对照组的小白鼠饲养在正常环境，一段时间后统计每只小白鼠体重的增加量（单位：\( g \)）.

\begin{enumerate}
\item 设 \( X \) 表示指定的两只小白鼠中分配到对照组的只数，求 \( X \) 的分布列和数学期望；
\item 实验结果如下：

对照组的小白鼠体重的增加量从小到大排序为：

15.2,\allowbreak 18.8,\allowbreak 20.2,\allowbreak 21.3,\allowbreak 22.5,\allowbreak 23.2,\allowbreak 25.8,\allowbreak 26.5,\allowbreak 27.5,\allowbreak 30.1,\allowbreak

32.6,\allowbreak 34.3,\allowbreak 34.8,\allowbreak 35.6,\allowbreak 35.6,\allowbreak 35.8,\allowbreak 36.2,\allowbreak 37.3,\allowbreak 40.5,\allowbreak 43.2

实验组的小白鼠体重的增加量从小到大排序为：

7.8,\allowbreak 9.2,\allowbreak 11.4,\allowbreak 12.4,\allowbreak 13.2,\allowbreak 15.5,\allowbreak 16.5,\allowbreak 18.0,\allowbreak 18.8,\allowbreak 19.2,\allowbreak

19.8,\allowbreak 20.2,\allowbreak 21.6,\allowbreak 22.8,\allowbreak 23.6,\allowbreak 23.9,\allowbreak 25.1,\allowbreak 28.2,\allowbreak 32.3,\allowbreak 36.5

\begin{enumerate}
\item 求 40 只小鼠体重的增加量的中位数 \( m \)，再分别统计两样本中小于 \( m \) 与不小于 \( m \) 的数据的个数，完成如下列联表；
\item 根据（1）中的列联表，能否有 95\% 的把握认为小白鼠在高浓度臭氧环境中与正常环境中体重的增加量有差异.
\end{enumerate}
\end{enumerate}

\begin{center}
\begin{tabular}{|c|c|c|}
\hline
& \( <m \) & \( \ge m \) \\
\hline
\text{对照组} &  &  \\
\hline
\text{实验组} &  &  \\
\hline
\end{tabular}
\end{center}

附：\( K^2=\dfrac{n\left( ad-bc \right)^2}{\left( a+b \right)\left( c+d \right)\left( a+c \right)\left( b+d \right)} \)，

\begin{center}
\begin{tabular}{c|ccc}
\hline
\( P\left( K^2\ge k \right) \) & 0.100 & 0.050 & 0.010 \\
\hline
\( k \) & 2.706 & 3.841 & 6.635 \\
\hline
\end{tabular}
\end{center}

\end{problem}

\begin{solution}
（1）依题意，\( X \) 的可能取值为 \( 0,1,2 \).

\(P(X=0)=\frac{C_{20}^0C_{20}^2}{C_{40}^2}=\frac{19}{78}\)，\(P(X=1)=\frac{C_{20}^1C_{20}^1}{C_{40}^2}=\frac{20}{39}\)，\(P(X=2)=\frac{C_{20}^2C_{20}^0}{C_{40}^2}=\frac{19}{78}\).

所以 \( X \) 的分布列为
\[
\begin{array}{c|ccc}
X & 0 & 1 & 2\\
\hline
P & \frac{19}{78} & \frac{20}{39} & \frac{19}{78}
\end{array}
\]

故
\(E(X)=0\times \frac{19}{78}+1\times \frac{20}{39}+2\times \frac{19}{78}=1\).

（2）（1）依题意，这 40 只小白鼠体重增加量的中位数是将两组数据合在一起，从小到大排后第 20 位与第 21 位数据的平均数. 观察数据可得第 20 位为 23.2，第 21 位数据为 23.6，

所以
\(m=\frac{23.2+23.6}{2}=23.4\).

故列联表为
\[
\begin{array}{c|ccc}
& <m & \ge m & \text{合计}\\
\hline
\text{对照组} & 6 & 14 & 20\\
\text{实验组} & 14 & 6 & 20\\
\hline
\text{合计} & 20 & 20 & 40
\end{array}
\]

（2）由（1）可得
\(K^2=\frac{40\times\left( 6\times 6-14\times 14 \right)^2}{20\times 20\times 20\times 20}=6.400>3.841\).

所以能有 95\% 的把握认为小白鼠在高浓度臭氧环境中与正常环境中体重的增加量有差异.
\end{solution}

\begin{problem}
已知直线 \( x-2y+1=0 \) 与抛物线 \( C:y^2=2px \)（\( p>0 \)）交于 \( A,B \) 两点，且 \( \lvert AB \rvert=4\sqrt{15} \).

\begin{enumerate}
\item 求 \( p \)；
\item 设 \( F \) 为 \( C \) 的焦点，\( M,N \) 为 \( C \) 上两点，\( \overrightarrow{FM}\cdot\overrightarrow{FN}=0 \)，求 \( \triangle MFN \) 面积的最小值.
\end{enumerate}

\end{problem}

\begin{solution}
（1）设 \( A(x_A,y_A) \)，\( B(x_B,y_B) \).

由
\(\begin{cases}
x-2y+1=0,\\
y^2=2px
\end{cases}\)
可得
\(y^2-4py+2p=0\)，
所以
\(y_A+y_B=4p\)，\(y_Ay_B=2p\).

于是
\[
\lvert AB \rvert=\sqrt{\left( x_A-x_B \right)^2+\left( y_A-y_B \right)^2}=\sqrt{5}\,\lvert y_A-y_B \rvert=\sqrt{5}\sqrt{\left( y_A+y_B \right)^2-4y_Ay_B}=4\sqrt{15}
\]

即
\(2p^2-p-6=0\).
因为 \( p>0 \)，解得
\(p=2\).

（2）因为 \( p=2 \)，所以抛物线方程为
\(y^2=4x\)，
其焦点为 \( F(1,0) \).

若直线 \( MN \) 的斜率为 0，可设 \( MN:y=t \). 代入 \( y^2=4x \) 得
\(x=\frac{t^2}{4}\)，
此时直线 \( MN \) 与抛物线只有 1 个公共点 \( \left( \frac{t^2}{4},t \right) \)，这与 \( M,N \) 为两个不同交点矛盾. 所以直线 \( MN \) 的斜率不可能为 0. 设直线 \( MN:x=my+n \)，\( M(x_1,y_1) \)，\( N(x_2,y_2) \).

由
\(\begin{cases}
y^2=4x,\\
x=my+n
\end{cases}\)
可得
\(y^2-4my-4n=0\)，
所以
\(y_1+y_2=4m\)，\(y_1y_2=-4n\).

又判别式
\(\Delta=16m^2+16n>0\)，
故
\(m^2+n>0\).

因为 \( \overrightarrow{FM}\cdot\overrightarrow{FN}=0 \)，所以
\(\left( x_1-1 \right)\left( x_2-1 \right)+y_1y_2=0\).

即
\(\left( my_1+n-1 \right)\left( my_2+n-1 \right)+y_1y_2=0\)，
也即
\(\left( m^2+1 \right)y_1y_2+m\left( n-1 \right)\left( y_1+y_2 \right)+\left( n-1 \right)^2=0\).

将 \( y_1+y_2=4m \)，\( y_1y_2=-4n \) 代入得 \(4m^2=n^2-6n+1\)，\(4\left( m^2+n \right)=\left( n-1 \right)^2>0\).

所以 \( n\neq 1 \)，且
\(n^2-6n+1\ge 0\)，
解得
\(n\ge 3+2\sqrt{2}\)或\(n\le 3-2\sqrt{2}\).

设点 \( F \) 到直线 \( MN \) 的距离为 \( d \)，则
\(d=\frac{\lvert n-1 \rvert}{\sqrt{1+m^2}}\).

又
\[
\lvert MN \rvert=\sqrt{\left( x_1-x_2 \right)^2+\left( y_1-y_2 \right)^2}=\sqrt{1+m^2}\,\lvert y_1-y_2 \rvert=\sqrt{1+m^2}\sqrt{16m^2+16n}
\]

由 \( 4m^2=n^2-6n+1 \) 得
\[
\lvert MN \rvert=\sqrt{1+m^2}\sqrt{4\left( n^2-6n+1 \right)+16n}=2\sqrt{1+m^2}\,\lvert n-1 \rvert
\]

所以
\[
S_{\triangle MFN}=\frac12\lvert MN \rvert d=\frac12\times 2\sqrt{1+m^2}\,\lvert n-1 \rvert\times \frac{\lvert n-1 \rvert}{\sqrt{1+m^2}}=\left( n-1 \right)^2
\]

而 \( n\ge 3+2\sqrt{2} \) 或 \( n\le 3-2\sqrt{2} \)，

所以当 \( n=3-2\sqrt{2} \) 时，
\(S_{\triangle MFN,\min}=\left( 2-2\sqrt{2} \right)^2=12-8\sqrt{2}\).

故 \( \triangle MFN \) 面积的最小值为 \( 12-8\sqrt{2} \).
\end{solution}

\begin{problem}
已知函数 \( f(x)=ax-\frac{\sin x}{\cos^3x} \)，\( x\in\left( 0,\frac{\pi}{2} \right) \).

\begin{enumerate}
\item 当 \( a=8 \) 时，讨论 \( f(x) \) 的单调性；
\item 若 \( f(x)<\sin2x \) 恒成立，求 \( a \) 的取值范围.
\end{enumerate}

\end{problem}

\begin{solution}
（1）
\[
f'(x)=a-\frac{\cos x\cos^3x+3\sin x\cos^2x\sin x}{\cos^6x}=a-\frac{\cos^2x+3\sin^2x}{\cos^4x}=a-\frac{3-2\cos^2x}{\cos^4x}
\]

令
\(t=\cos^2x\)，
则 \( t\in(0,1) \)，

于是
\(f'(x)=\frac{at^2+2t-3}{t^2}\).

当 \( a=8 \) 时，
\(f'(x)=\frac{8t^2+2t-3}{t^2}=\frac{\left( 2t-1 \right)\left( 4t+3 \right)}{t^2}\).

当 \( t\in\left( 0,\frac12 \right) \) 时，即 \( x\in\left( \frac{\pi}{4},\frac{\pi}{2} \right) \) 时，\( f'(x)<0 \)；

当 \( t\in\left( \frac12,1 \right) \) 时，即 \( x\in\left( 0,\frac{\pi}{4} \right) \) 时，\( f'(x)>0 \).

所以 \( f(x) \) 在 \( \left( 0,\frac{\pi}{4} \right) \) 上单调递增，在 \( \left( \frac{\pi}{4},\frac{\pi}{2} \right) \) 上单调递减.

（2）设
\(g(x)=f(x)-\sin2x\).

则
\[
g'(x)=f'(x)-2\cos2x=\frac{at^2+2t-3}{t^2}-2\left( 2t-1 \right)=a+2-4t+\frac2t-\frac3{t^2}
\]

记
\(\varphi(t)=a+2-4t+\frac2t-\frac3{t^2}\).

则
\[
\varphi'(t)=-4-\frac2{t^2}+\frac6{t^3}=-\frac{2\left( t-1 \right)\left( 2t^2+2t+3 \right)}{t^3}>0
\]

所以
\(\varphi(t)<\varphi(1)=a-3\).

\circled{1} 若 \( a\le 3 \)，则
\(g'(x)=\varphi(t)<a-3\le 0\).

即 \( g(x) \) 在 \( \left( 0,\frac{\pi}{2} \right) \) 上单调递减，所以
\(g(x)<g(0)=0\).

因此当 \( a\in\left( -\infty,3 \right] \) 时，\( f(x)<\sin2x \) 符合题意.

\circled{2} 若 \( a>3 \)，当 \( t\to 0 \) 时，\(\frac2t-\frac3{t^2}=-3\left( \frac1t-\frac13 \right)^2+\frac13\to -\infty\)，所以 \( \varphi(t)\to -\infty \)，而 \( \varphi(1)=a-3>0 \).

故存在 \( t_0\in(0,1) \) 使得 \( \varphi(t_0)=0 \). 设 \( x_0\in\left( 0,\frac{\pi}{2} \right) \) 满足 \( \cos^2x_0=t_0 \).

当 \( t\in(t_0,1) \) 时，\( \varphi(t)>0 \)，对应当 \( x\in(0,x_0) \) 时，\( g'(x)>0 \)，所以 \( g(x) \) 单调递增.

于是当 \( x\in(0,x_0) \) 时，\(g(x)>g(0)=0\)，不合题意.

综上，\( a \) 的取值范围为 \( \left( -\infty,3 \right] \).
\end{solution}

\begin{problem}
已知点 \( P(2,1) \)，直线 \( l \) 的参数方程为 \( x=2+t\cos\alpha \)，\( y=1+t\sin\alpha \)（\( t \) 为参数），\( \alpha \) 为 \( l \) 的倾斜角，\( l \) 与 \( x \) 轴正半轴、\( y \) 轴正半轴分别交于 \( A,B \) 两点，且 \( \lvert PA \rvert\cdot \lvert PB \rvert=4 \).

\begin{enumerate}
\item 求 \( \alpha \)；
\item 以坐标原点为极点，\( x \) 轴正半轴为极轴建立极坐标系，求 \( l \) 的极坐标方程.
\end{enumerate}

\end{problem}

\begin{solution}
（1）因为 \( l \) 与 \( x \) 轴、\( y \) 轴正半轴交于 \( A,B \) 两点，所以
\(\frac{\pi}{2}<\alpha<\pi\).

令 \( x=0 \)，得
\(t_1=-\frac{2}{\cos\alpha}\).
令 \( y=0 \)，得
\(t_2=-\frac{1}{\sin\alpha}\).

所以 \( \lvert PA \rvert\cdot \lvert PB \rvert=\lvert t_1t_2 \rvert=\left\lvert \frac{2}{\sin\alpha\cos\alpha} \right\rvert=\left\lvert \frac{4}{\sin2\alpha} \right\rvert=4 \).

故 \(\sin2\alpha=\pm 1\).

于是 \(2\alpha=\frac{\pi}{2}+k\pi\).
解得
\(\alpha=\frac{\pi}{4}+\frac{k\pi}{2}\)（\(k\in \Z\)）.

又 \( \frac{\pi}{2}<\alpha<\pi \)，所以
\(\alpha=\frac{3\pi}{4}\).

（2）由（1）可知，直线 \( l \) 的斜率为
\(\tan\alpha=-1\)，
且过点 \( (2,1) \)，

所以直线 \( l \) 的普通方程为
\(y-1=-\left( x-2 \right)\)，
即
\(x+y-3=0\).

由 \(x=\rho\cos\theta\)，\(y=\rho\sin\theta\) 可得直线 \( l \) 的极坐标方程为 \(\rho\cos\theta+\rho\sin\theta-3=0\).
\end{solution}

\begin{problem}
设 \( a>0 \)，函数 \( f(x)=2\lvert x-a \rvert-a \).

\begin{enumerate}
\item 求不等式 \( f(x)<x \) 的解集；
\item 若曲线 \( y=f(x) \) 与 \( x \) 轴所围成的图形面积为 2，求 \( a \).
\end{enumerate}

\end{problem}

\begin{solution}
（1）若 \( x\le a \)，则 \(f(x)=2a-2x-a=a-2x\).

由 \( f(x)<x \) 得
\(a-2x<x\)，
即
\(3x>a\)，
解得
\(\frac{a}{3}<x\le a\).

若 \( x>a \)，则
\(f(x)=2x-2a-a=2x-3a\).

由 \( f(x)<x \) 得
\(2x-3a<x\)，
即
\(x<3a\).

故
\(a<x<3a\).

综上，不等式的解集为
\(\left( \frac{a}{3},3a \right)\).

（2）
\[
f(x)=
\begin{cases}
-2x+a,& x\le a,\\
2x-3a,& x>a.
\end{cases}
\]

画出 \( f(x) \) 的草图，则 \( f(x) \) 与 \( x \) 轴围成 \( \triangle ABC \).

其中高为 \( a \)，两交点分别为
\(A\left( \frac{a}{2},0 \right)\)，\(B\left( \frac{3a}{2},0 \right)\)，
所以
\(\lvert AB \rvert=a\).

于是
\(S_{\triangle ABC}=\frac12\lvert AB \rvert\cdot a=\frac12a^2=2\).

解得 \(a=2\).
\end{solution}
